\section{Sturm--Liouville Problems}

\subsection{Abstract Eigenvalue Problems}

Recall from \textit{IA Vectors and Matrices}, a \term{linear map} \(\functiondomain{A}{V_N}{V_N}\) is called \term{Hermitian} if \(\vb{x} \vdot A \vb{y} = \pqty{A \vb{x}} \vdot \vb{y}\) for all \(\vb{x}, \vb{y} \in V_N\), where we use the complex inner product given by
\begin{equation}
    \vb{x} \vdot \vb{y} = \sum_{i=1}^N x_i \conjugate{y_i} = \vb{x}\transpose \vb{y}\conjstar
\end{equation}
This is equivalent, in matrix notation, saying that \(A = A\hermitian\). Hermitian maps have the following properties.

\begin{enumerate}[label=(\roman*)]
    \item Eigenvalues of Hermitian maps are real.
    \item Eigenvectors of Hermitian maps corresponding to distinct eigenvalues are orthogonal.
    \item There exists a set \(\Bqty{\vb{y}_n}_{n = 1}^{N}\) of eigenvectors such that every \(\vb{x} \in V_N\) can be written as
    \begin{subequations}
        \begin{equation}
            \vb{x} = \sum_{n = 1}^{N} \hat{x}_n \vb{y}_n
        \end{equation}
        where
        \begin{equation}
            \hat{x}_n = \frac{\vb{x} \vdot \vb{y}_n}{\vb{y}_n \vdot \vb{y}_n} = \frac{\vb{x} \vdot \vb{y}_n}{\norm{\vb{y}_n}^2}.
        \end{equation}
    \end{subequations}
\end{enumerate}

Our question is, does this generalise in the following way?

Let \(V\) be a vector space of functions \(\functiondomain{f}{\ccintv{a}{b}}{\RR}\) (\(V\) will be a subspace of \(\contClass{2}{\ccintv{a}{b}}\)). Equip \(V\) with the \term{inner product}
\begin{equation}
    \inner[w]{f}{g} = \int_{a}^{b} f\pqty{x} \conjugate{g\pqty{x}} w\pqty{x} \dd{x},
\end{equation}
where \(w\pqty{x}\) is real and \(w > 0\) on \(\oointv{a}{b}\), referred to as the \term{weight function}. Define the \term{norm}
\begin{equation}
    \norm[w]{f} = \sqrt{\inner[w]{f}{f}},
\end{equation}
and we write
\begin{equation}
    \inner[1]{f}{g} = \inner{f}{g}.
\end{equation}

\begin{definition}[Self-Adjoint Operator]
    A linear operator \(L\) on \(V\) is called \term{self-adjoint} on \(\pqty{V, \inner[w]{\cdot}{\cdot}}\), if
    \begin{equation}
        \inner[w]{L y_1}{y_2} = \inner[w]{y_1}{L y_2}
    \end{equation}
    for all \(y_1, y_2 \in V\).
\end{definition}

We say \(\pqty{\lambda, y} \in \CC \times \pqty{V \setminus \Bqty{0}}\) is an \term{eigenvalue-eigenfunction pair} for \(L\), if \(Ly = \lambda_y\).

\begin{proposition}[Properties of Self-Adjoint Operators]
    If \(L\) is self-adjoint on \(\pqty{V, \inner[w]{\cdot}{\cdot}}\), then:
    \begin{enumerate}[label=(\roman*)]
        \item eigenvalues of \(L\) are real;
        \item eigenfunctions of \(L\) corresponding to distinct eigenvalues are orthogonal with respect to \(\inner[w]{\cdot}{\cdot}\);
        \item there exists a set \(\Bqty{y_n}_{n = 1}^{\infty}\) of eigenfunctions such that every \(f \in V\) can be written as
        \begin{equation}
            f = \sum_{n = 1}^{\infty} \hat{f}_n y_n,
        \end{equation}
        where
        \begin{equation}
            \hat{f}_n = \frac{\inner[w]{f}{y_n}}{\inner[w]{y_n}{y_n}} = \frac{\inner[w]{f}{y_n}}{\norm[w]{y_n}^2}.
        \end{equation}
    \end{enumerate}
\end{proposition}

\begin{proof}
    \begin{enumerate}[label=(\roman*)]
        \item Suppose \(L y = \lambda y\) for \(y \neq 0\). Then
        \begin{subequations}
            \begin{align}
                \pqty{\lambda - \conjugate{\lambda}} \norm[w]{y}^2 &= \pqty{\lambda - \conjugate{\lambda}} \inner[w]{y}{y} \\
                &= \inner[w]{\lambda y}{y} - \inner[w]{y}{\lambda y} \\
                &= \inner[w]{Ly}{y} - \inner[w]{y}{Ly} \\
                &= 0,
            \end{align}
        \end{subequations}
        since \(L\) is self-adjoint. So \(\lambda = \conjugate{\lambda}\), i.e. \(\lambda \in \RR\).

        \item Suppose \(L y_1 = \lambda_1 y_1\) and \(L y_2 = \lambda_2 y_2\) for \(y_1, y_2 \neq 0\) and \(\lambda_1 \neq \lambda_2\). Then
        \begin{subequations}
            \begin{align}
                \pqty{\lambda_1 - \lambda_2} \inner[w]{y_1}{y_2} &= \inner[w]{\lambda_1 y_1}{y_2} - \inner[w]{y_1}{\lambda_2 y_2} \\
                &= \inner[w]{Ly_1}{y_2} - \inner[w]{y_1}{Ly_2} \\
                &= 0,
            \end{align}
        \end{subequations}
        since \(L\) is self-adjoint. So \(\inner[w]{y_1}{y_2} = 0\).

        \item This proof is left as an exercise to the reader. In particular, the famous mathematician David Hilbert says it is true, so it must be true.
    \end{enumerate}
\end{proof}

\small
Essentially, at the end of \textit{II Linear Analysis} we may prove that
\begin{equation}
    \norm[w]{f - \sum_{n = 1}^{N} \hat{f}_n y_n} \to 0
\end{equation}
as \(n \to \infty\).
\normalsize

\subsection{Self-Adjointness and Boundary Conditions}

We will study problems of the form
\begin{equation}
    Ly = \lambda_y \qfor a < x < b
\end{equation}
and with some supplementary boundary conditions at \(x = a\) and \(x = b\). Here, we restrict \(L\) to be a special type of differential operator.

\begin{definition}[Sturm--Liouville Operator]
    A \term{Sturm--Liouville operator} is a second-order differential operator of the form
    \begin{equation}
        L y = \frac{1}{w} \bqty{- \dv{x}  \pqty{p \dv{x}} y + q y},
    \end{equation}
    where \(p, q\) and \(w\) are real-valued functions, and \(p\) and \(w\) are positive on \(\oointv{a}{b}\). We call \(w\) the \term{weight function associated with \(L\)}.
\end{definition}

Rearranging gives the common form of a Sturm--Liouville problem:
\begin{equation}
    - \dv{x} \bqty{p \dv{y}{x}} + qy = \lambda w y
\end{equation}
on \(a < x < b\).

We will impose boundary conditions by forcing \(y\) to belong to some vector space \(V\).

\begin{definition}[Singular and Non-Singular Endpoints]
    Given a Sturm--Liouville operator on \(\oointv{a}{b}\), we say an endpoint \(c \in \Bqty{a, b}\) is \term{singular} if \(p\pqty{c} = 0\), and \term{non-singular} otherwise.
\end{definition}

We will only specify boundary conditions at non-singular endpoints. At non-singular endpoints, we will impose real, homogeneous boundary conditions of the form
\begin{equation}
    \alpha_c y\pqty{c} + \beta_c y'\pqty{c} = 0,
\end{equation}
where \(\alpha_c, \beta_c \in \RR\) and \(\alpha^2_c + \beta^2_c \neq 0\). Note that if \(y_1, y_2\) satisfy boundary conditions of this form, then so does \(c_1 y_1 + c_2 y_2\) for constants \(c_1, c_2\), i.e., we got a vector space structure.

So given Sturm--Liouville operator, we work on a vector space of the form
\begin{equation}
    V = \set{y \in \contClass{2}\pqty{\ccintv{a}{b}}}{y \text{ satisfies real, homogeneous boundary conditions at non-singular endpoints}}
\end{equation}

We equip \(V\) with an inner product \(\inner[w]{\cdot}{\cdot}\), where \(w\) is the weight function associated with \(L\). 

\begin{example}
    Consider the Sturm--Liouville problem
    \begin{equation}
        - \dv{x} \bqty{\cos \pqty{\frac{x}{2}} \dv{y}{x}} + \sin \pqty{\frac{x}{2}} y = \lambda y
    \end{equation}
    where \(p = \cos \pqty{\flatfrac{x}{2}},  q = \sin \pqty{\flatfrac{x}{2}}, w = 1\). A suitable vector space would be
    \begin{equation}
        V = \set{y \in \contClass{2}\pqty{\ccintv{0}{\pi}}}{y\pqty{0} = 0}.
    \end{equation}

    Note that \(p\pqty{\pi} = 0\), \(p\pqty{0} \neq 0\) so only the endpoint \(0\) is non-singular. The inner product is therefore given by
    \begin{equation}
        \innerproduct{f}{g} = \int_{0}^{\pi} f\pqty{x} \conjugate{g\pqty{x}} \dd{x}.
    \end{equation}
\end{example}

\begin{proposition}
    Let \(V\) be a vector space associated with a Sturm--Liouville operator \(L\) on \(\oointv{a}{b}\). Then \(L\) is self-adjoint on \(\pqty{V, \inner[w]{\cdot}{\cdot}}\), where \(w\) is a weight function corresponding to \(L\).
\end{proposition}

\begin{proof}
    Take \(y_1, y_2 \in \contClass{2} \pqty{\ccintv{a}{b}}\). Then we have
    \begin{subequations}
        \label{eq:sturm-ioiville-self-adjointness}%
        \begin{align}
            \inner[w]{L y_1}{y_2} - \inner[w]{y_1}{L y_2} &= \int_{a}^{b} \frac{1}{w} \bqty{- \pqty{p y_1'}' + q y_1} \conjugate{y_2} w \dd{x} - \int_{a}^{b} y_1 \frac{1}{w} \bqty{- \pqty{p \conjugate{y_2}'}' + q \conjugate{y_2}} w \dd{x} \\
            &= \int_{a}^{b} \bqty{y_1 \pqty{p \conjugate{y_2}'}' - \pqty{p y_1'}' \conjugate{y_2}} \dd{x} \\
            &= \int_{a}^{b} \bqty{p \pqty{y_1 \conjugate{y_2}' - y_1' \conjugate{y_2}}}' \dd{x} \\
            &= \bqty{p \pqty{x} W\pqty{y_1, \conjugate{y_2}}}_{x = a}^{b}, \label{eq:sturm-liouville-wronskian}%
        \end{align}
    \end{subequations}
    where \(W\pqty{x}\) is the \term{Wronskian} of \(y_1\) and \(y_2\) (recall from \textit{IA Differential Equations}).

    Now assume \(y_1, y_2 \in V\). If \(c \in \Bqty{a, b}\) is singular, then \(p \pqty{c} = 0\), i.e., the right-hand side of \autoref{eq:sturm-liouville-wronskian} vanishes at \(x = c\).

    If \(c \in \Bqty{a, b}\) is non-singular, they behave boundary conditions of the form
    \begin{equation}
        \alpha_c y\pqty{c} + \beta_c y'\pqty{c} = 0,
    \end{equation}
    where \(\alpha_c, \beta_c \in \RR\) and \(\alpha_c^2 + \beta_c^2 \neq 0\).Since \(\alpha_c\) and \(\beta_c\) are real, \(\conjugate{y}\) also solves this equation. So we have
    \begin{equation}
        \pmqty{y_1\pqty{c} & y_1'\pqty{c} \\ \conjugate{y_2}\pqty{c} & \conjugate{y_2}'\pqty{c}} \pmqty{\alpha_c \\ \beta_c} = 0.
    \end{equation}
    at this non-singular endpoint. Since \(\alpha_c^2 + \beta_c^2 \neq 0\), the vector \(\pmqty{\alpha_c \\ \beta_c}\) is non-zero, and so the determinant of the matrix, which is exactly \(W\pqty{y_1, \conjugate{y_2}}\pqty{c}\), must vanish. So at a non-singular endpoint, the right-hand side of \autoref{eq:sturm-liouville-wronskian} also vanishes.

    This shows that the left-hand side of \autoref{eq:sturm-ioiville-self-adjointness} vanishes, i.e., \(L\) is self-adjoint on \(\pqty{V, \inner[w]{\cdot}{\cdot}}\).
\end{proof}